\subsection{子空间与张成集}
\index{subspace|(}
在\nearbyexample{PlaneThruOriginSubsp}中，我们见到了一个向量空间，它是 $\Re^2$ 的一个子集：一条过原点的直线。
在那里，向量空间 $\Re^2$ 的内部还包含着另一个
向量空间，即这条直线。

\begin{definition} \label{df:Subspace}
%<*df:Subspace>
对于任意向量空间，
\definend{子空间}\index{vector space!subspace}\index{subspace!definition}
是指在继承的运算之下，其本身也是一个向量空间的子集。
%</df:Subspace>
\end{definition}

\begin{example}  \label{ex:PlaneSubspRThree}
这个过原点的平面
\begin{equation*}
  P=\set{\colvec{x \\ y \\ z}\suchthat x+y+z=0}
\end{equation*}
是 \( \Re^3 \) 的一个子空间。
正如定义所要求的，
该平面的运算继承自更大的空间，
也就是说，
向量在 $P$ 中的加法与在 $\Re^3$ 中的加法一样
\begin{equation*}
   \colvec{x_1 \\ y_1 \\ z_1}+\colvec{x_2 \\ y_2 \\ z_2}
   =\colvec{x_1+x_2 \\ y_1+y_2 \\ z_1+z_2}
\end{equation*}
且标量乘法也与 $\Re^3$ 中的相同。
为证明 $P$ 是子空间，我们只需注意它是一个子集，
然后验证它是一个空间。
我们不必检验全部十个条件，只检验两个封闭性条件。
关于对加法封闭，注意若
两个加项分别满足
$x_1+y_1+z_1=0$ 与 $x_2+y_2+z_2=0$，则其和满足
$(x_1+x_2)+(y_1+y_2)+(z_1+z_2)=(x_1+y_1+z_1)+(x_2+y_2+z_2)=0$。
关于对标量乘法封闭，若
$x+y+z=0$，则其标量倍满足
$rx+ry+rz=r(x+y+z)=0$。
\end{example}

\begin{example}   \label{ex:SubspacesRTwo}
\( \Re^2 \) 中的 \( x \) 轴
是一个子空间，其中
加法与标量乘法运算都是
继承而来的那些运算。
\begin{equation*}
  \colvec{x_1 \\ 0}
    +
  \colvec{x_2 \\ 0}
    =
  \colvec{x_1+x_2 \\ 0}
  \qquad
  r\cdot\colvec{x \\ 0}
  =\colvec{rx \\ 0}
\end{equation*}
与前面的例题一样，要从定义直接验证它是子空间，
我们只需注意它是一个
子集，然后检验它满足
向量空间定义中的各个条件。
例如两个封闭性条件都
满足：第二个分量为零的两个向量相加，所得向量的第二个分量仍为零；
标量乘以第二个分量为零的向量，
所得向量的第二个分量仍为零。
\end{example}

\begin{example}
$\Re^2$ 的另一个子空间是
它的平凡子空间。
\begin{equation*}
  \set{\colvec[r]{0 \\ 0}}
\end{equation*}
\end{example}

%<*ProperSubspace>
任何向量空间都有一个平凡子空间
\( \set{\zero\,} \)。\index{subspace!trivial}%
\index{trivial subspace}
与之相反的极端是，任何向量空间都以自身作为自己的一个子空间。
不是整个空间的子空间称为
\definend{真}\index{subspace!proper}%
\index{proper subspace} 子空间。
%</ProperSubspace>

\begin{example}  \label{ex:LinSubspPolyThree}
不是 $\Re^n$ 的向量空间也有子空间。
三次多项式的空间
\( \set{a+bx+cx^2+dx^3\suchthat a,b,c,d\in\Re} \)
有一个由全部一次多项式
\( \set{m+nx\suchthat m,n\in\Re} \)
组成的子空间。
\end{example}

\begin{example}
另一个不是 $\Re^n$ 的子集的子空间的例子出现在向量空间的定义之后。
\nearbyexample{ex:RealValuedFcns}中全体一元实值函数构成的空间
\( \set{f\suchthat \map{f}{\Re}{\Re} } \) 含有\nearbyexample{ex:SpaceSatisfyingDiffQ}中满足
限制 $(d^2\,f/dx^2)+f=0$ 的那些函数组成的子空间。
\end{example}

\begin{example} \label{ex:OperNotInherit}
该定义要求
加法与标量乘法运算
必须是继承自更大空间的那两种运算。
集合 \( S=\set{1} \) 是 \( \Re^1 \) 的一个子集。
而且在运算 $1+1=1$ 与 $r\cdot 1=1$ 之下，
集合 $S$ 是一个向量空间，确切地说，是一个平凡空间。
然而 $S$ 不是 \( \Re^1 \) 的子空间，因为那些运算不是
继承的运算，毕竟 \( \Re^1 \) 中当然有 \( 1+1=2 \)。
\end{example}

\begin{example}  \label{cex:RPlusNotSubSp}
子空间本身是向量空间，因此必须满足封闭性
条件。
集合 \( \Re^+ \) 不是向量空间 \( \Re^1 \) 的子空间，
因为在继承的运算下它对标量乘法不封闭：
若 \( \vec{v}=1 \)，则 \( -1\cdot\vec{v}\not\in\Re^+ \)。
\end{example}

下面的结果表明，\nearbyexample{cex:RPlusNotSubSp} 是典型的。
一个子集若非空且使用继承的运算，
则它不成其为子空间的唯一可能就是不封闭。

\begin{lemma}     \label{th:SubspIffClosed} \index{subspace!closed}
%<*lm:SubspIffClosed>
对于向量空间的一个非空子集 \( S \)，在继承的
运算之下，下列诸命题
等价。\appendrefs{equivalence of statements}   %\spacefactor=1000
\begin{tfae}
  \item \( S \) 是该向量空间的一个子空间
  \item \( S \) 对两个向量的线性组合封闭：
    对任意向量 \( \vec{s}_1,\vec{s}_2\in S \) 与标量 \( r_1,r_2 \)，
    向量 \( r_1\vec{s}_1+r_2\vec{s}_2 \) 属于 \( S \)
  \item \( S \) 对任意多个向量的线性组合封闭：
    对任意向量 \( \vec{s}_1,\ldots,\vec{s}_n\in S \) 与标量
    \( r_1, \ldots,r_n \)，
    向量 \( r_1\vec{s}_1+\cdots+r_n\vec{s}_n \) 是 \( S \) 的元素。
\end{tfae}
%</lm:SubspIffClosed>
\end{lemma}
\noindent 简言之，一个子集是
子空间当且仅当它对线性组合封闭。

\begin{proof}
%<*pf:SubspIffClosed0>
“下列诸命题等价”是指每一对
命题都等价。
\begin{equation*}
  (1)\!\iff\!(2)
  \qquad
  (2)\!\iff\!(3)
  \qquad
  (3)\!\iff\!(1)
\end{equation*}
我们将通过确立
\( (1)\implies (3)\implies (2)\implies (1) \) 来证明这一等价性。
采取这一策略的理由是如下观察：蕴含关系
\( (1)\implies (3) \) 与 \( (3)\implies (2) \) 都很容易，于是我们只需
论证 \( (2)\implies (1) \)。
%</pf:SubspIffClosed0>

%<*pf:SubspIffClosed1>
设 \( S \) 是向量空间
$V$ 的一个非空子集，且对两个向量的组合封闭。
我们将通过检验各项条件来证明 $S$ 是一个向量空间。

向量空间的定义中有五条关于加法的条件。
第一，关于对加法封闭，若
\( \vec{s}_1,\vec{s}_2\in S \)，则 \( \vec{s}_1+\vec{s}_2\in S \)，
因为它是
两个向量的组合，
而我们正假设~\( S \) 对这种组合封闭。
第二，对任意 \( \vec{s}_1,\vec{s}_2\in S \)，由于加法
继承自 \( V \)，
\( S \) 中的和 \( \vec{s}_1+\vec{s}_2 \)
等于 \( V \) 中的和 \( \vec{s}_1+\vec{s}_2 \)，
而后者等于
\( V \) 中的和 \( \vec{s}_2+\vec{s}_1 \)（因为 $V$ 是向量空间，其加法可交换），
进而后者又等于 \( S \) 中的和 \( \vec{s}_2+\vec{s}_1 \)。
第三个条件的论证与第二个类似。
第四，考虑 \( V \) 的零向量，注意
$S$ 对两个向量的线性组合的封闭性给出：
\( 0\cdot\vec{s}+0\cdot\vec{s}=\zero \) 是 \( S\) 的元素
（其中 \( \vec{s} \) 是非空集 \( S \) 的任一成员）；
容易检验 \( \zero \) 在继承的运算下充当 \( S \) 的加法
单位元。
第五个条件满足，因为对任意 \( \vec{s}\in S \)，
对两个向量的线性组合的封闭性表明
\( 0\cdot\zero+(-1)\cdot\vec{s} \) 是~\( S \) 的元素，而且它显然是
\( \vec{s} \) 在继承的运算下的
加法逆元。
%</pf:SubspIffClosed1>
关于标量乘法诸条件的验证类似；
见\nearbyexercise{exer:SubspIffClosed}。
\end{proof}

以后我们通常通过检验一个子集
满足命题~(2) 来验证它是子空间。

\begin{remark}
在本章开头，我们把向量空间作为线性组合在其中“有意义”的集合加以介绍。
% The vector space definition has ten conditions but eight of them, the
% conditions not about closure, simply ensure that referring to the
% operations by the names `addition' and `scalar multiplication' is sensible.
% For instance, calling an operation `$+$' would not
% be odd unless it is commutative, as in condition~(2). 
% The proof above checks that the subspace satisfies these 
% eight non-closure conditions because they hold in the
% surrounding vector space, provided that the nonempty set $S$ satisfies
% \nearbylemma{th:SubspIffClosed}'s statement~(2) 
% (e.g., commutativity of addition in $S$ follows right from
% commutativity of addition in $V$).
% There is a second way in which we can regard vector spaces as structures in 
% which linear combinations are ``sensible.''
% It has to do with the closure conditions.
\nearbylemma{th:SubspIffClosed} 的诸命题 (1)-(3)
表明我们总能对形如
$r_1\vec{s}_1+r_2\vec{s}_2$
的表达式给出意义，
也就是说，所描述的向量在集合~$S$ 中。

作为对比，考虑由分量之和不小于零的二元向量组成的集合 $T$。
在这里我们不能随手写出任意的线性组合如 $2\vec{t}_1-3\vec{t}_2$
并确信其结果是 $T$ 的元素。
\end{remark}

\nearbylemma{th:SubspIffClosed} 提示我们，理解向量空间的一种好方式
是把它看作无限制线性组合的集合。
接下来的两个例子取一些空间，把对它们的描述改写成这种形式。

\begin{example}
要证明 $\Re^3$ 的这个过原点的平面子集
\begin{equation*}
  S=\set{\colvec{x \\ y \\ z}\suchthat x-2y+z=0}
\end{equation*}
在列向量通常的加法与标量乘法
运算下是子空间，我们可以检验它非空且对
两个向量的线性组合封闭。
但还有另一种方法。
把 $x-2y+z=0$ 看作只含一个方程的线性方程组并加以参数化：
将首变量用自由变量表示为 $x=2y-z$。
\begin{equation*}
     S
     =\set{\colvec{2y-z \\ y \\ z}\suchthat y,z\in\Re}
     =\set{y\colvec[r]{2 \\ 1 \\ 0}+
            z\colvec[r]{-1 \\ 0 \\ 1}\suchthat y,z\in\Re}
    \tag{$*$}
\end{equation*}
现在，为证明它是子空间，考虑
$r_1\vec{s}_1+r_2\vec{s}_2$。
每个 $\vec{s}_i$ 都是 ($*$) 中那两个向量的线性组合，
所以这是一个线性组合的线性组合。
\begin{equation*}
  r_1\cdot(y_1\colvec[r]{2 \\ 1 \\ 0}+
            z_1\colvec[r]{-1 \\ 0 \\ 1})
  +
  r_2\cdot(y_2\colvec[r]{2 \\ 1 \\ 0}+
            z_2\colvec[r]{-1 \\ 0 \\ 1})
\end{equation*}
线性组合引理（Lemma~One.III.\ref{lm:LinearCombinationLemma}）
表明总和是那两个向量的线性组合，
于是 \nearbylemma{th:SubspIffClosed} 的命题~(2) 得到满足。
\end{example}

\begin{example} \label{ex:ParamSubspace}
这是全体 \( \nbyn{2} \) 矩阵的空间 $\matspace_{\nbyn{2}}$ 的一个子空间。
\begin{equation*}
  L=\set{\begin{mat}
         a  &0  \\
         b  &c
       \end{mat}
       \suchthat a+b+c=0}
\end{equation*}
为参数化，把条件写成 $a=-b-c$。
\begin{equation*}
  L
  =\set{\begin{mat}
         -b-c  &0  \\
         b     &c
       \end{mat}
       \suchthat b,c\in\Re}
  =\set{b\begin{mat}[r]
         -1    &0  \\
         1     &0
       \end{mat}
       +c\begin{mat}[r]
         -1    &0  \\
         0     &1
       \end{mat}
       \suchthat b,c\in\Re}
\end{equation*}
如上所述，我们已把这个子空间描述为无限制线性
组合的集合。
为证明它是子空间，注意来自
$L$ 的向量的线性组合是线性组合的线性组合，因此命题~(2)
成立。
\end{example}

% Parametrization is easy, but important.
% We shall use it often.

\begin{definition} \label{df:Span}
%<*df:Span>
向量空间的一个非空子集 \( S \) 的\definend{张成}\index{span}\index{sets!span of}\index{closure}（或称
\definend{线性闭包}）是由 \( S \) 中向量的一切线性组合组成的集合。
\begin{equation*}
  \spanof{S} =\{ c_1\vec{s}_1+\cdots+c_n\vec{s}_n
            \suchthat c_1,\ldots, c_n\in\Re
            \text{\ 且\ } \vec{s}_1,\ldots,\vec{s}_n\in S \}
\end{equation*}
向量空间的空子集的张成是它的平凡子空间。
%</df:Span>
\end{definition}

%<*NotationForSpan>
\noindent 关于张成的记号没有完全标准的。
这里用的方括号很常见，
“$\mbox{span}(S)$” 与 “$\mbox{sp}(S)$” 也很常用。
%</NotationForSpan>

\begin{remark}
在第一章，当我们证明了可以把齐次线性方程组的解
集写成
$\set{c_1\vec{\beta}_1+\cdots+c_k\vec{\beta}_k\suchthat
  c_1,\ldots,c_k\in\Re}$
之后，我们把它描述为由诸 $\smash{\vec{\beta}}$ “生成”的集合。
现在我们称它为
$\set{\vec{\beta}_1,\ldots,\vec{\beta}_k}$ 的张成。

也应从那个证明中回忆起：
空集的张成定义为集合 \( \set{\zero} \)，
这是出于这样一个约定：平凡的线性组合，即零个向量的组合，其和为~\( \zero \)。
此外，把空集的张成定义为平凡子空间
也很方便，因为它使得像下一个结果那样的定理
无需为空集另立例外。
\end{remark}

\begin{lemma}   \label{le:SpanIsASubsp}
%<*lm:SpanIsASubsp>
在向量空间中，任何子集的张成都是子空间。
%</lm:SpanIsASubsp>
\end{lemma}

\begin{proof}
%<*pf:SpanIsASubsp>
若子集 \( S \)
为空，则按定义其张成是平凡
子空间。
若 \( S\) 非空，则由\nearbylemma{th:SubspIffClosed}，我们
只需检验张成 \( \spanof{S} \) 对两个元素的线性组合封闭。
取该张成中的一对向量，
\( \vec{v}=c_1\vec{s}_1+\cdots+c_n\vec{s}_n \) 与
\( \vec{w}=c_{n+1}\vec{s}_{n+1}+\cdots+c_m\vec{s}_m \)，
则线性组合
\begin{multline*}
  p\cdot(c_1\vec{s}_1+\cdots+c_n\vec{s}_n)+
       r\cdot(c_{n+1}\vec{s}_{n+1}+\cdots+c_m\vec{s}_m)  \\
  =
  pc_1\vec{s}_1+\cdots+pc_n\vec{s}_n
    +rc_{n+1}\vec{s}_{n+1}+\cdots+rc_m\vec{s}_m
\end{multline*}
是~\( S \) 中元素的线性组合，
因此是 \( \spanof{S} \) 的元素
（来自 $\vec{v}$ 的某些 $\vec{s}_i$ 可能与来自 $\vec{w}$ 的某些 $\vec{s}_j$ 相同，但这无关紧要）。
%</pf:SpanIsASubsp>
\end{proof}

这个引理的逆
也成立：任何子空间都是某个集合的张成，
因为子空间显然是它自身的张成，即由其全体成员
构成的集合的张成。
于是，向量空间的一个子集是子空间当且仅当它是一个张成。
这符合下述
直觉：理解向量空间的好方式是把它看作
线性组合在其中有意义的集合。

合起来，\nearbylemma{th:SubspIffClosed} 与
\nearbylemma{le:SpanIsASubsp} 表明：向量空间的子集 $S$ 的
张成是包含 $S$ 的全体成员的最小子空间。

\begin{example}   \label{ex:SpanSingVec}
在任意向量空间 \( V \) 中，对任意向量 \( \vec{v}\in V \)，集合
\( \set{r\cdot\vec{v} \suchthat r\in\Re} \) 是 \( V \) 的子空间。
例如，对任意向量 \( \vec{v}\in\Re^3 \)，
包含该向量的过原点直线
\( \set{k\vec{v}\suchthat k\in\Re } \) 是 \( \Re^3 \) 的子空间。
即使 $\vec{v}$ 是零向量，这也是对的，此时
它是退化直线，即平凡子空间。
\end{example}

\begin{example}
这个集合的张成
是整个 $\Re^2$。
\begin{equation*}
  \set{\colvec[r]{1 \\ 1},\colvec[r]{1 \\ -1}}
\end{equation*}
我们知道该张成是~$\Re^2$ 的某个子空间。
要检验它就是整个~$\Re^2$，
必须证明 $\Re^2$ 的任何成员都是这两个向量的
线性组合。
于是我们问：~对以实数 $x$ 与~$y$ 为分量的哪些向量，
存在标量 $c_1$ 与 $c_2$ 使下式成立？
\begin{equation*}
   c_1\colvec[r]{1 \\ 1}+c_2\colvec[r]{1 \\ -1}=\colvec{x \\ y}
  \tag{$*$}
\end{equation*}
高斯消元法
\begin{equation*}
  \begin{linsys}{2}
    c_1  &+  &c_2  &=  &x  \\
    c_1  &-  &c_2  &=  &y
  \end{linsys}
  \grstep{-\rho_1+\rho_2}
  \begin{linsys}{2}
    c_1  &+  &c_2    &=  &x\hfill  \\
         &   &-2c_2  &=  &-x+y
  \end{linsys}
\end{equation*}
经回代给出 $c_2=(x-y)/2$ 与 $c_1=(x+y)/2$。
这表明对任意 $x,y$ 都存在
适当的系数 $c_1,c_2$ 使~($*$) 成立\Dash 我们可以把 $\Re^2$ 的任意元素写成这两个
给定向量的线性组合。
例如，取 $x=1$ 与 $y=2$，系数 $c_2=-1/2$ 与
$c_1=3/2$ 即可。
\end{example}

既然张成是子空间，而我们又知道
理解子空间的好方式是
将其描述参数化，那么我们可以尝试用这种方式来理解一个集合的张成。

\begin{example}
在次数不超过二的多项式空间~\( \polyspace_2 \) 中，考虑
集合 \( S=\set{3x-x^2, 2x} \) 的张成。
由张成的定义，它是这两个向量的无限制线性组合
$\set{c_1(3x-x^2)+c_2(2x)\suchthat c_1,c_2\in\Re}$ 组成的集合。
显然这个张成中的多项式的常数项必为零。
这个必要条件也是充分的吗？

我们问的是：~对 $\polyspace_2$ 的哪些成员 $a_2x^2+a_1x+a_0$
存在 $c_1$ 与 $c_2$ 使
$a_2x^2+a_1x+a_0=c_1(3x-x^2)+c_2(2x)$？
系数相等的多项式相等，因此
我们要找 $a_2$、$a_1$ 与 $a_0$ 上的条件，
使这个三元组成为下面方程组的解。
\begin{equation*}
  \begin{linsys}{2}
    -c_1  &   &     &=  &a_2   \\
    3c_1  &+  &2c_2 &=  &a_1   \\
          &   &0    &=  &a_0
  \end{linsys}
\end{equation*}
高斯消元法与回代给出
$c_1=-a_2$，以及 $c_2=(3/2)a_2+(1/2)a_1$，以及 $0=a_0$。
于是
%the only condition on elements  $a_0+a_1x+a_2x^2$ of the span
%is the condition that we knew:
只要没有常数项，即 $a_0=0$，
我们就能给出系数 $c_1$ 与 $c_2$，
把该多项式描述为张成中的一个元素。
例如，对多项式 $0-4x+3x^2$，系数
$c_1=-3$ 与 $c_2=5/2$ 即可。
所以给定集合的张成是
$\spanof{S}=\set{a_1x+a_2x^2\suchthat a_1,a_2\in\Re}$。

顺带地，这表明
集合 \( \set{x,x^2} \) 张成同一个子空间。
一个空间可以有不止一个张成集。
张成这个子空间的另外两个集合是
\( \set{x,x^2,-x+2x^2} \) 与
\( \set{x,x+x^2,x+2x^2,\ldots\,} \)。
% (Usually we prefer to work with spanning sets that have only
% a small number of members.)
\end{example}

\begin{example}  \label{ex:SubspRThree}
下图展示了我们目前已知的 \( \Re^3 \) 的子空间：平凡
子空间、过原点的直线、
过原点的平面，以及整个空间。
（当然，图示只画出了无穷多种情形中的少数几个。
线段把子集与其
超集连接起来。）
在下一节我们将证明 $\Re^3$ 没有其他
类型的子空间，因此实际上这张图列出了全部子空间。

这把每个子空间
描述为成员个数最少的集合的张成。
据此，各子空间
自然分成一些层次\Dash 平面在同一层，
直线在另一层，
等等。
\begin{center}
  \setlength{\unitlength}{4pt}
  \begin{picture}(75,38)(0,-2) %subspaces of R-three
      \thinlines
      \put(45,31){\makebox(0,0)[bl]{
                        \scriptsize \( %\Re^3=
                                   \set{x\colvec[r]{1 \\ 0 \\ 0}
                                              +y\colvec[r]{0 \\ 1 \\ 0}
                                              +z\colvec[r]{0 \\ 0 \\ 1}} \)} }
    %Next the dimension two subspaces
      \put(43,31){\line(-4,-1){25} } % connects R3 to 2,a
      %set 2,a:
      \put(0,22){\makebox(0,0)[l]{\scriptsize\( \set{x\colvec[r]{1 \\ 0 \\ 0}
                                                 +y\colvec[r]{0 \\ 1 \\ 0} }\) }}
      \put(48,29){\line(-3,-1){12} } % connects R3 to 2,b
      %set 2,b:
      \put(20,22){\makebox(0,0)[l]{\scriptsize\( \set{x\colvec[r]{1 \\ 0 \\ 0}
                                                 +z\colvec[r]{0 \\ 0 \\ 1} }\) }}
      \put(53,29){\line(-1,-1){3}} % connects R3 to 2,c
      %set 2,c:
      \put(40,22){\makebox(0,0)[l]{\scriptsize\( \set{x\colvec[r]{1 \\ 1 \\ 0}
                                                 +z\colvec[r]{0 \\ 0 \\ 1} }\) }}
      \put(60,22){\makebox(0,0)[l]{$\cdots$} }
    %Next the dimension one subspaces
      \put(7,18){\line(-1,-4){1} } % connects 2,a to 1,a
      \put(22,18){\line(-3,-1){13} } % connects 2,c to 1,a
      %set 1,a:
      \put(0,10){\makebox(0,0)[l]{\scriptsize\( \set{x\colvec[r]{1 \\ 0 \\ 0}} \)} }
      \put(12,18){\line(1,-2){2} } % connects 2,a to 1,b
      %set 1,b:
      \put(12,10){\makebox(0,0)[l]{\scriptsize\( \set{y\colvec[r]{0 \\ 1 \\ 0}} \)} }
      \put(14,18){\line(2,-1){10} } % connects 2,a to 1,c
      %set 1,c:
      \put(24,10){\makebox(0,0)[l]{\scriptsize\( \set{y\colvec[r]{2 \\ 1 \\ 0}} \)} }
      \put(46,18){\line(-1,-1){3.25} } %connects 2,c to 1,d
      %set 1,d:
      \put(36,10){\makebox(0,0)[l]{\scriptsize\( \set{y\colvec[r]{1 \\ 1 \\ 1}} \)} }
      \put(48,10){\makebox(0,0)[l]{$\cdots$} }
    %Finally, the trivial subspace.
      \put(9,7.5){\line(4,-1){30} } %connects 1,a to trivial
      \put(18.9,7.7){\line(3,-1){20} } %connects 1,b to trivial
      \put(30.4,6.8){\line(2,-1){10} } %connects 1,c to trivial
      \put(41,6){\line(1,-2){1.5} } %connects 1,d to trivial
      \put(45,0){\makebox(0,0){\scriptsize\( \set{\colvec[r]{0 \\ 0 \\ 0} } \)} }
  \end{picture}
\end{center}
\end{example}

到本章目前为止我们已经看到：要研究
线性组合的性质，合适的框架是
对这些组合封闭的集合。
在第一小节中我们引入了这类集合，即向量空间，
并见到了丰富多样的例子。
在本小节中我们又见到了
更多的空间，它们是其他空间的子空间。
在所有这些多样性中有一个共同点。
上面的\nearbyexample{ex:SubspRThree}
把它揭示了出来：~向量空间与子空间最好被理解为张成，
尤其是少量向量的张成。
下一节将研究具有极小性的张成集。






